\begingroup\small
\begin{longtable}{l p{7.2cm} l}
\hline
statement & title & type \\
\hline
\endhead
Lemma~II.2.1 & powerful numbers & proved \\
Proposition~II.2.2 & Walker {\cite[Sections~2 and~3]{Walker1976}} & literature \\
Proposition~II.2.3 & Walker {\cite[Theorems 2.2, 2.5, 3.2 and 3.5]{Walker1976}} & literature \\
Proposition~II.3.1 & the known pairs & computed \\
Proposition~II.3.2 & completeness inside the families & computed \\
Proposition~II.4.1 & explicit lower bound & computed \\
Definition~II.5.1 & comparison function; a heuristic & heuristic \\
Proposition~II.5.2 & the families with a square overtake the comparison function & computed \\
Question~II.5.3 & convergence of the rates & open \\
Proposition~II.6.1 & the middles that are not divisible by four & computed \\
Proposition~II.6.2 & the count is a sum over towers & proved \\
Corollary~II.6.3 & lower bound at every height & computed \\
Proposition~II.6.4 & localization of the pair count & proved \\
Corollary~II.6.5 & what the localization says, and what it does not say & proved \\
Remark~II.6.6 & where Tao's argument is weakest & proved \\
Question~II.6.7 & growth of the tower rate & open \\
Proposition~II.7.1 & pairs on one Pell equation & proved \\
Proposition~II.7.2 & the third number & computed \\
Question~II.8.1 & the kernels that divide $U_1$ & open \\
\hline
\end{longtable}
\endgroup

The types: \emph{proved}, proved in this part; the proof may use published results (6); \emph{literature}, due to the cited source, sometimes with a proof given here for completeness (2); \emph{computed}, the result of an exact computation, or a proof that uses one of our computations (7); \emph{heuristic}, a heuristic, used as a yardstick only (1); \emph{open}, an open question (3).
